\chapter*{Contribuciones personales}
\label{cap:contribucionesPersonales}
\addcontentsline{toc}{chapter}{Contribuciones personales}

El presente Trabajo de Fin de Grado ha sido desarrollado de forma conjunta por los dos autores. Ambos estudiantes han participado en la definición del problema, en las reuniones de seguimiento con los tutores, en la toma de decisiones técnicas, en la revisión de resultados y en la elaboración de la memoria. No obstante, dentro de este trabajo común, cada integrante ha asumido responsabilidades predominantes en distintas fases del proyecto, que se detallan a continuación.

\section*{Beatriz Bueno Almagro}

Mi contribución al proyecto se ha centrado principalmente en la definición técnica de los modelos, el análisis de la arquitectura empleada, el ajuste experimental de la salida del sistema, la integración de técnicas de explicabilidad y la redacción y revisión de distintas partes de la memoria. Desde las primeras fases del trabajo participé en la comprensión del problema planteado por la ICTS BioImagen Complutense, especialmente en la traducción de las necesidades del usuario final a requisitos técnicos concretos: segmentar automáticamente el cerebro y la lesión isquémica, calcular volúmenes y generar una salida interpretable dentro de un flujo de trabajo biomédico asistido.

Durante la fase de investigación inicial colaboré en la revisión del estado del arte, prestando especial atención a las técnicas de visión por computador aplicadas a imagen biomédica, a las arquitecturas de segmentación médica, al uso de U-Net y al aprendizaje por transferencia en contextos con conjuntos de datos reducidos. Esta investigación fue necesaria para justificar decisiones metodológicas importantes, como el uso de una arquitectura U-Net con encoder preentrenado frente a alternativas más complejas, y para argumentar la necesidad de incorporar técnicas de Inteligencia Artificial Explicable en un sistema aplicado a imagen médica.

Una parte relevante de mi trabajo consistió en estudiar y documentar la arquitectura de los modelos de segmentación. Participé en la descripción de la U-Net 2D empleada, en la justificación del uso de un encoder ResNet50 preentrenado con RadImageNet y en la explicación de las decisiones de implementación asociadas al decoder, las conexiones de salto, la salida sigmoidal y el planteamiento de la segmentación como problema binario. También contribuí a relacionar estas decisiones con las características del proyecto: tamaño reducido del conjunto de datos, naturaleza biomédica de las imágenes, necesidad de segmentación precisa y limitaciones computacionales propias del alcance de un TFG.

Asimismo, trabajé de forma destacada en el ajuste de la salida del modelo, especialmente en el análisis del umbral de binarización aplicado sobre la predicción probabilística. En el caso del modelo de isquemia, este ajuste resultó especialmente importante debido a la tendencia del sistema a sobresegmentar regiones de la imagen y generar falsos positivos. Por ello, participé en la comparación de distintos valores de umbral y en la interpretación de sus efectos sobre las métricas de evaluación y sobre la coherencia visual de las máscaras obtenidas. Esta parte del trabajo permitió justificar por qué no basta con aceptar directamente la salida probabilística del modelo, sino que es necesario adaptar el umbral al comportamiento concreto de la tarea.

Otra de mis principales aportaciones fue la integración y análisis de técnicas de explicabilidad. Trabajé en el estudio de métodos como Grad-CAM, Grad-CAM++ e Integrated Gradients, así como en su adaptación conceptual a un problema de segmentación, donde la salida del modelo no es una única clase sino una máscara espacial. Esta tarea incluyó reflexionar sobre qué significa explicar una predicción en segmentación médica y cómo deben interpretarse los mapas generados. En la memoria se plantea esta explicabilidad no como una validación clínica directa, sino como una herramienta para analizar si el modelo basa sus predicciones en regiones anatómica o patológicamente coherentes.

También me encargué especialmente de la interpretación de los resultados de explicabilidad. Esto incluyó el análisis cualitativo de los mapas obtenidos para los modelos, la comparación entre las regiones activadas por los métodos de XAI y las máscaras de segmentación, y la discusión de las diferencias entre métodos. En particular, participé en la interpretación de por qué Grad-CAM y Grad-CAM++ tienden a producir explicaciones más localizadas y visualmente interpretables, mientras que Integrated Gradients puede mostrar patrones más finos pero también más dispersos. Esta discusión permitió complementar las métricas de segmentación con una lectura más crítica del comportamiento interno del modelo.

En relación con la memoria, participé en la redacción de apartados técnicos vinculados a la arquitectura, el aprendizaje por transferencia, la estrategia de entrenamiento, el ajuste del umbral y la explicabilidad. Además, contribuí a la revisión global del documento, corrigiendo repeticiones, mejorando la cohesión entre capítulos y ajustando el tono académico. Esta revisión fue especialmente importante para que la memoria no quedase como una simple descripción de código, sino como una explicación razonada de las decisiones tomadas durante el desarrollo.

También participé en la revisión de la trazabilidad entre la implementación y la memoria, comprobando que las decisiones descritas en el documento se correspondieran con los notebooks desarrollados y con los resultados obtenidos. Esta revisión incluyó la comprobación de la coherencia entre metodología, figuras, tablas, métricas y discusión, evitando que el texto presentase afirmaciones no respaldadas por la implementación o por los experimentos realizados.

Por último, contribuí a la preparación de la defensa del trabajo en los apartados relacionados con la motivación técnica del sistema, la justificación de la arquitectura, la elección del umbral, la explicabilidad y las limitaciones del modelo. También participé en la revisión final de figuras, tablas, referencias y formulaciones técnicas, comprobando que las afirmaciones incluidas en la memoria estuvieran alineadas con la implementación real y con los resultados obtenidos.

\section*{Leonardo Sáez Biffi}

Mi contribución al proyecto se ha centrado principalmente en la implementación técnica del sistema, el procesamiento de los datos, el entrenamiento de los modelos de segmentación, la construcción del flujo de ejecución en notebooks y la redacción de varias partes fundamentales de la memoria. Desde las primeras fases del trabajo participé en el análisis del problema planteado por la ICTS BioImagen Complutense y en la definición de una solución basada en segmentación automática de imágenes de resonancia magnética preclínica.

Una de mis principales aportaciones fue la preparación del conjunto de datos para el entrenamiento de los modelos. Esta tarea incluyó la lectura y organización de los casos de resonancia magnética, la correspondencia entre imágenes y máscaras, la preparación de las segmentaciones de cerebro e isquemia y la adaptación de los datos al formato requerido por la red. También participé en la resolución de problemas derivados de la estructura real del dataset, como la existencia de múltiples máscaras asociadas a un mismo corte o la necesidad de mantener una correspondencia correcta entre las imágenes originales y las regiones de interés anotadas manualmente.

Durante el desarrollo del sistema implementé buena parte del pipeline de preprocesamiento de imágenes. Este proceso incluyó la lectura de volúmenes de resonancia, la extracción de cortes bidimensionales, la normalización de intensidades, el redimensionado de las imágenes al tamaño de entrada del modelo y la generación de tensores adecuados para el entrenamiento. También trabajé en la adaptación de imágenes monocanales a una entrada de tres canales, necesaria para emplear un encoder preentrenado dentro de la arquitectura U-Net. Estas transformaciones fueron fundamentales para asegurar que los datos pudieran introducirse de forma consistente en los modelos de aprendizaje profundo.

Otra aportación central fue la implementación y entrenamiento de los modelos de segmentación. Participé en el desarrollo de los notebooks de entrenamiento tanto para la segmentación del cerebro como para la segmentación de la lesión isquémica, utilizando una arquitectura U-Net con encoder ResNet50 y pesos preentrenados de RadImageNet. En esta fase se trabajó con estrategias de \textit{transfer learning}, congelación inicial del encoder y posterior ajuste fino de capas, así como con funciones de pérdida adecuadas para segmentación binaria, especialmente en un contexto de fuerte desbalance entre fondo y lesión.

Además, desarrollé parte del flujo de inferencia del sistema, integrando los dos modelos dentro de un mismo pipeline. En este flujo, el modelo de cerebro se utiliza primero para delimitar la región anatómica principal y, posteriormente, el modelo de isquemia realiza la segmentación de la lesión sobre la imagen procesada. Esta integración permitió acercar el sistema al uso final previsto: cargar una resonancia, procesar sus cortes, generar máscaras de segmentación y obtener resultados cuantitativos asociados. También participé en la preparación de notebooks de prueba que permitieran recorrer este proceso de forma secuencial y reproducible.

Otra parte relevante de mi trabajo fue la implementación del cálculo de volúmenes a partir de las máscaras generadas. Para ello se tuvo en cuenta la información espacial de las imágenes originales, de forma que el volumen segmentado pudiera expresarse en unidades físicas y no únicamente como número de píxeles. Esta funcionalidad resulta especialmente importante en el contexto biomédico del proyecto, ya que el objetivo de la herramienta no es solo producir una máscara visual, sino facilitar una cuantificación útil de la estructura cerebral y de la lesión. En este sentido, la implementación del cálculo volumétrico conecta directamente la salida del modelo con una necesidad práctica del usuario final.

También participé en la adaptación de los notebooks para que pudieran ser ejecutados por terceros de forma más sencilla. Esto incluyó la limpieza del código, la organización de las celdas, la incorporación de mensajes explicativos, la gestión básica de errores y la preparación de un flujo de uso más accesible en Google Colab. Esta parte del trabajo fue importante para que el sistema no quedase únicamente como un conjunto de scripts de desarrollo, sino como una herramienta reproducible y evaluable por usuarios externos al proceso de implementación.

En cuanto a la memoria, participé de forma directa en la redacción de la introducción, donde se presenta la motivación del proyecto, el contexto de colaboración con ICTS BioImagen Complutense y la necesidad de automatizar tareas de segmentación y cuantificación en imágenes de resonancia magnética. También redacté el primer apartado de metodología, centrado en la descripción del conjunto de datos y su organización, explicando la estructura de los casos, las imágenes utilizadas, las máscaras disponibles y la relación entre los datos originales y los datos empleados para el entrenamiento.

Además, me encargué de redactar gran parte del capítulo de resultados, incluyendo la presentación de las métricas obtenidas, la comparación entre las segmentaciones automáticas y las máscaras de referencia, y el análisis del comportamiento de los modelos en los distintos casos evaluados. Esta parte del trabajo fue especialmente importante para conectar la implementación con una evaluación objetiva del sistema, mostrando no solo el funcionamiento del pipeline, sino también sus aciertos, errores y limitaciones. También participé en la preparación de tablas, gráficas y ejemplos visuales empleados para apoyar esta discusión.

Asimismo, participé en la redacción de las conclusiones y del trabajo futuro. En este apartado contribuí a sintetizar el grado de cumplimiento de los objetivos planteados, las principales aportaciones del sistema desarrollado y las limitaciones que permanecen abiertas. También colaboré en la identificación de posibles líneas futuras, como la ampliación del conjunto de datos, la mejora del modelo de lesión, la validación con más casos y la evolución del prototipo hacia una herramienta más completa y usable en un entorno biomédico real.

Durante la fase de evaluación, colaboré en la obtención de resultados cuantitativos y cualitativos. Esto incluyó el cálculo de métricas de segmentación, la comparación entre predicciones y máscaras de referencia, y la generación de ejemplos visuales que permitieran analizar el comportamiento de los modelos en cortes concretos. Estos resultados fueron posteriormente utilizados para construir las tablas y figuras incluidas en la memoria, así como para apoyar la discusión sobre los puntos fuertes y limitaciones del sistema.

Finalmente, participé en la revisión técnica de la memoria, asegurando que la descripción del código, del entrenamiento y del funcionamiento del pipeline coincidiera con la implementación real desarrollada durante el proyecto. También contribuí a la preparación de la defensa en los apartados más relacionados con la estructura del dataset, el preprocesamiento de datos, el flujo de ejecución, el entrenamiento de los modelos, la obtención de resultados y el cálculo de volúmenes.